\documentclass[
    newtxsf,
    % mtpro2,
    twoside,openany,
    % hideanswer
]{litesolution}
\usepackage{weyounghead}
% 中文字体设置，很好看的一款字体，可以去GitHub上搜索安装后使用
\setCJKmainfont{LXGW WenKai}[AutoFakeBold=3,ItalicFont=*]

% \CenterWallPaper{.5}{WeYoung}
\fancyhead[OL]{\itshape\color{DarkGreen}\href{http://weixin.qq.com/r/DBIABcHEqgnWKIVmb0eI}{\faIcon{weixin}~微信公众号：未央学社}}
\fancyhead[ER]{\color{MidnightBlue}\href{https://qm.qq.com/q/OlkSiy9W04}{\faIcon{qq}~\textit{未央数学营：}782790567}}

\begin{document}

% \chapterimage{\UPFadingImage[trim={0 {.9\paperheight} 0 0},clip]{cover2}}
\chapter*{2024年\textsf{HDU}「概率论与数理统计」期中试题解答}
\weyounghead{理学院}{A0714040}{待定ing}{幻想大白}{https://qm.qq.com/q/OlkSiy9W04}{未央数学营}

\section{单项选择题（每题3分，共30分）}

\begin{problem}
    设A和B为任意两个事件。则下列等式一定成立的是
    \begin{tasks}(2)
        \task $P(\overline A\ \overline B)=P(\overline A)P(\overline B)$
        \task $P(\overline A\ \overline B) = P(\overline A) P(\overline B)$
        \task $P(\overline AB)=P(B)-P(AB)$
        \task $P(A\cup B) = P(A) + P(B)$
    \end{tasks}
\end{problem}
\begin{solution}

\end{solution}

\begin{problem}
    设随机变量$x$的概率密度函数为$f(x)$。且$f(-x)=f(x),F(x)$是$X$的分布函数，则对于任意实数$a$有
    \begin{tasks}(2)
        \task$F(-a)=1-\int_{0}^{a}f(x)dx$
        \task$F(-a)=1/2-\int_{0}^{a}f(x)dx$
        \task$F(-a)=F(a)$
        \task$F(-a)=2F(a)-1$
    \end{tasks}
\end{problem}
\begin{solution}

\end{solution}

\begin{problem}
    设随机变量$X\sim N(2,4)$。则$P(X<1)$的值为
    \begin{tasks}(4)
        \task$\phi(2)$
        \task$\phi(4)$
        \task$\phi(0.5)$
        \task$\phi(-0.5)$
    \end{tasks}
\end{problem}
\begin{solution}

\end{solution}

\begin{problem}
    已知函数$F(x)=
    \begin{cases}
        0,      & x<0\\
        \sin x, & 0\leq x<\pi,\\
        1,      & x\geq\pi
    \end{cases}$.
    下列关于$F(x)$说法正确的是
    \begin{tasks}(2)
        \task 不是某随机变量的分布函数
        \task 是某随机变量的分布函数
        \task 是某离散型随机变量的分布函数
        \task 是某连续型随机变量的分布函数
    \end{tasks}
\end{problem}
\begin{solution}

\end{solution}

\begin{problem}
    设测机变量$X$的指率密度为$f_X(x)$，$x\in R$,令$Y=-2X$，则$Y$的概率密度$f_Y(y)$为
    \begin{tasks}(4)
        \task $2f_X(2y)$
        \task $-2f_X\ab(-\frac y2)$
        \task $\frac12f_X\ab(-\frac y2)$
        \task $-\frac12f_X\ab(-\frac y2)$
    \end{tasks}
\end{problem}
\begin{solution}

\end{solution}


\begin{problem}
    已知随机变量$X\sim b(n,p)$，且$E(X)=2$，$D(X)=1.6$，则参数$n$，$p$的值为
\begin{tasks}(4)
    \task $n=4$，$p=0.5$
    \task $n=10$，$p=0.2$
    \task $n=10$，$p=0.8$
    \task $n=5$，$p=0.4$
\end{tasks}

\end{problem}
\begin{solution}

\end{solution}

\begin{problem}
    设X和Y为两随机变量,且$P\{X\geq0,\ Y\geq0\}=3/7$，$P\{X\geq0\}=P\{Y\geq0\}=4/7$，则$P\{\max(X,Y)\geq0\}=$
    \begin{tasks}(4)
        \task $1/7$
        \task $2/7$
        \task $3/7$
        \task $5/7$
    \end{tasks}
\end{problem}
\begin{solution}

\end{solution}

\begin{problem}
    已知随机变量$X$的密度函数为$f(x)=
    \begin{cases}
        \e^{-x/2}/2, & x>0 \\
        0,           & \text{其他}
    \end{cases}$，另$Y=2X+\e^{-2X}$，则$E(Y)$的值为
    \begin{tasks}(4)
        \task $21/5$
        \task $5$
        \task $3/4$
        \task $4/3$
    \end{tasks}
\end{problem}
\begin{solution}

\end{solution}

\begin{problem}
    设$P(AB)=0$，则
    \begin{tasks}(2)
        \task $A$和$B$互不相容
        \task $A$和$B$相互独立
        \task $P(A)=0$或$P(B)=0$
        \task $P(A-B)=P(A)$
    \end{tasks}
\end{problem}
\begin{solution}

\end{solution}

\begin{problem}
    设随机变量$X$的分布函数为$F(x)=
    \begin{cases}
        0,   & x<-1\\
        0.4, & -1\leq x<1\\
        0.8, & 1\leq x<3\\
        1,   & x\geq3
    \end{cases}$，则$P\{X<1|X\neq1\}=$
    \begin{tasks}(4)
        \task $2/3$
        \task $2/5$
        \task $3/5$
        \task $4/5$
    \end{tasks}
\end{problem}
\begin{solution}

\end{solution}

\section{填空题（共5题，每题3分,共15分）}

\begin{problem}
    从不同的5双袜子中任取4只,问这4只袜子中至少有两只配成一双的概率为\ans{}
\end{problem}
\begin{solution}

\end{solution}

\begin{problem}
    设A，B，C是三个事件，$P(A)=\frac12$，$P(B)=\frac13$，$P(C)=\frac15$，$P(AB)=P(BC)=0$，$P(AC)=1/6,则P(\overline A\ \overline{BC})=$\ans{}
\end{problem}
\begin{solution}

\end{solution}

\begin{problem}
    设随机变量$X\sim\pi(\lambda),\lambda>0,且P\left\{X=2\right\}=2P\{X=1\}$,则$\lambda=$\ans{}
\end{problem}
\begin{solution}

\end{solution}

\begin{problem}
    随机变量X的概率密度为$f(x)=
    \begin{cases}
        2x, &0<x<1\\
        0,  &\text{其它}
    \end{cases}$，$Y$表示对$X$三次独立重复观察中事件$\{X\leq1/3\}$出现的次数，求$P\{Y=2\}=$\ans{}（结果用分数表示）.
\end{problem}
\begin{solution}

\end{solution}

\begin{problem}
    设随机变量$X\sim N(0,4)$，$Y\sim U(-1,1)$,且X与Y相互独立，则$D(X-3Y+1)=$\ans{}.
\end{problem}
\begin{solution}

\end{solution}

\section{计算题（共55分）}

\begin{problem}[本题8分]
    一篮球运动员的投篮命中率为$45\%$. 以$X$表示他首次投中时累计已投篮的次数，求
    \begin{tasks}[label=\arabic*.](2)
        \task 写出$X$的分布律.
        \task 计算$X$取偶数的概率.
    \end{tasks}
\end{problem}
\begin{solution}

\end{solution}

\begin{problem}[本题15分]
    设随机变量:(X,Y)的概率密度函数
    \[f(x,y)=
    \begin{cases}
        1, & |y|<x,0<x<1\\
        0, & \text{else}
    \end{cases}\]
    \begin{tasks}[label=\arabic*.](2)
        \task 求$f_X(x)$，$f_Y(y)$.
        \task $X$与$Y$是否相互独立？说明理由.
        \task 求条件密度 $f_{X|Y}(x|1/2)$
        \task 计算$D(X)$的值（建议画图）.
    \end{tasks}
\end{problem}
\begin{solution}

\end{solution}

\begin{problem}[本题15分]
    设随机变量$X$与$Y$相互独立,下表是二维随机变量$(X,Y)$的分布律及边缘分布律中的部分数值
    \begin{table}[htbp]
        \centering
        \begin{tabular}{ccccc}
            \toprule
            $X$\textbackslash$Y$ & $-2$ & $0$ & $1$ & $P\{X=x_i\}=p_i$\\
            \midrule
            $1$ & (a) & $1/8$ & (b) & (c) \\
            $2$ & $1/8$ & (d) & (e) & (f) \\
            $P\{Y=y_j\}=p_j$ & $1/6$ & (g) & (h) & $1$\\
            \bottomrule
        \end{tabular}
    \end{table}
    \begin{tasks}[label=\arabic*.](2)
        \task! 填写表中空白(a)至(h)填入正确的数值，并写出计算过程.
        \task 关于$XY$的分布律.
        \task $\text{Cov}(X,-r)$
    \end{tasks}
\end{problem}
\begin{solution}

\end{solution}

\begin{problem}
    设随机变量X的概率密度函数为$f(x)=
    \begin{cases}
        x,   & 0\leq x<1\\
        a-x, & 1\leq x<2\\
        0,   & \text{其它}
    \end{cases}$.
    \begin{tasks}[label=\arabic*.](3)
        \task 求$a$的值.
        \task 求$X$的分布函数$F(x)$.
        \task 求P\{$1/2<X<3/2$\}.
    \end{tasks}
\end{problem}
\begin{solution}\leavevmode
    \begin{enumerate}
        \item 
        \begin{enumerate*}
            \item a
            \item b
            \item c
            \item d
            \item e
            \item f
            \item g
            \item h
        \end{enumerate*}
    \end{enumerate}
\end{solution}

\begin{problem}[本题7分]
    设随机变量$Z=(aX+3Y)^2$，$E(X)=E(Y)=0$，$D(X)=4$，$D(Y)=16$，$\rho_{XY}=-0.5$. 问当$a$取何值时，$E(Z)$的值最小，并求$E(Z)$的最小值
\end{problem}
\begin{solution}

\end{solution}

\end{document}
